\chapter{The Workflow Model}
\label{ch:workflow-model}

\newthought{The workflow is step two} of the per-point loop: it turns a parameter point into
observables. This short chapter explains the shared ideas behind both workflow backends---how
modules connect, what order things run in, and the flowchart \Jarvis\ draws---before the next
chapters dive into each backend.

\section{Two backends, one job}
\label{sec:wf-backends}

A workflow is built from \emph{modules}. Each module takes some inputs and produces some named
observables. \Jarvis\ offers two kinds of modules, and a card uses one kind or the other:

\begin{description}[style=nextline, leftmargin=1.2em]
  \item[\yamlkey{Calculators}] wrap external programs --- you give \Jarvis\ the commands to run
        and the files to read and write (Chapter~\ref{ch:calculators}).
  \item[\yamlkey{Operas}] call in-process Python operators --- no files, no shell, just a
        function call (Chapter~\ref{ch:operas}).
\end{description}

\begin{jtip}
Reach for \yamlkey{Operas} when the calculation already exists as an operator (it is the fastest
path), and \yamlkey{Calculators} when you must drive a compiled \textsc{hep} package through
files. The rest of this part covers both.
\end{jtip}

\section{Modules and their order}
\label{sec:wf-order}

Modules rarely run in isolation---one often needs another's output. You declare those
dependencies with \yamlkey{required\_modules}, and \Jarvis\ uses them to build a dependency
graph and run modules in a valid order:

\begin{yamlwin}
Modules:
  - name: Spectrum
    required_modules: []          # runs first
  - name: DarkMatter
    required_modules: [Spectrum]  # waits for Spectrum
\end{yamlwin}

\noindent Here \yamlkey{DarkMatter} only starts once \yamlkey{Spectrum} has produced its
observables. A module with \yamlkey{required\_modules: []} has no prerequisites and can run as
soon as the parameter point is ready.

\begin{jnote}
The scan parameters themselves are provided by a built-in \yamlkey{Parameters} layer that runs
first. A module that consumes scan variables directly can declare
\yamlkey{required\_modules: [Parameters]} so it is scheduled after that layer.
\end{jnote}

\section{The order inside a calculator}
\label{sec:wf-calc-order}

Within a single calculator module, the execution order is fixed and always the same:

\begin{enumerate}
  \item \textbf{write input files} --- turn the parameter point into the program's input format;
  \item \textbf{run the external commands} --- execute the program;
  \item \textbf{read output files} --- extract the observables you asked for.
\end{enumerate}

\noindent Keeping this order in mind makes calculator cards (Chapter~\ref{ch:calculators}) easy
to read: every one is just ``write, run, read.''

\section{The workflow flowchart}
\label{sec:wf-flowchart}

When a run starts, \Jarvis\ writes a \emph{semantic} description of your workflow to
\file{images/<scan>/flowchart.json} --- the modules, their inputs and outputs, and the
dependency edges between them. If \JPlot\ is installed, \Jarvis\ also asks it to render a
picture, \file{images/<scan>/flowchart.png}.

\begin{jnote}
\Jarvis\ always writes \file{flowchart.json}; the \textsc{png} is only produced when \JPlot\ is
available. Pass \cmd{-{}-skip-draw-flowchart} to skip the picture while still exporting the
\textsc{json} (Chapter~\ref{ch:cli}).
\end{jnote}

\noindent The flowchart is a quick way to confirm \Jarvis\ understood your workflow: check that
the modules appear in the order you expect and that the dependency arrows match your
\yamlkey{required\_modules}.

\section{Where to go next}
\label{sec:wf-next}

\begin{itemize}
  \item External programs: Chapter~\ref{ch:calculators}, then the I/O formats in
        Chapter~\ref{ch:io-types}.
  \item In-process operators: Chapter~\ref{ch:operas}.
  \item Shared libraries that must be built before any of this: Chapter~\ref{ch:libdeps}.
\end{itemize}
